Acierto & «Ahora bien, el otro gran eje del reporte no pasa por las favoritas actuales, sino por las acciones que podrían recibir más dinero si MSCI habilita una futura reclasificación.» & POS & POS & La oración identifica una oportunidad de potencial entrada de capitales ('recibir más dinero') para ciertas acciones, lo cual constituye una señal favorable en un… \\
\addlinespace
Acierto & «Mientras que, a nivel regional, fue el Noreste la que arrojó un número más preocupante (39\%).» & NEG & NEG & El uso del adjetivo 'preocupante' para calificar el índice en la región Noreste introduce una carga valorativa negativa sobre la situación socioeconómica de esa zona. \\
\addlinespace
Acierto & «Ese punto sigue siendo importante porque muchas veces se confunde el valor nominal con el dinero que efectivamente entra a la cuenta todos los meses.» & NEU & NEU & La oración es una observación analítica y aclaratoria sobre la distinción entre valor nominal y poder adquisitivo real, sin expresar una valoración positiva o negativa… \\
\addlinespace
Acierto (era MIXTO) & «Aunque Höli Haus no brinda financiación directa, Tizado aclaró que existen bancos dispuestos a otorgar créditos al comprador final, en condiciones que dependen del perfil de cada cliente.» & MIXTO \flecha{} NEU & NEU & La oración presenta una señal negativa (la falta de financiación directa por parte de la empresa) compensada por una señal positiva (la disponibilidad de alternativas de… \\
\addlinespace
Error & «En enero, se amplió el tope de US\$ 24.000 anuales de ese entonces, que, a su vez, antes era de US\$ 12.000.» & POS & NEU & La ampliación de un límite o tope para la entrada de divisas representa una medida favorable para los profesionales del sector dentro del marco de formalización descrito. \\
\addlinespace
Error & «Puesto a explicar hacia dónde va Aleph, Taratuta señala que los consumidores siempre llegan antes que las empresas y que su desafío es esa anticipación.» & NEU & NEG & La oración es una declaración descriptiva sobre la visión estratégica y el desafío operativo de la empresa, sin presentar una valoración positiva o negativa sobre el… \\
